\documentclass[11pt]{article}
% ======================================================================
%  examstyle.tex — EPFL-style exam layout for the weekly Time Series exams
%  Shared by every Exam_Week_NN(.tex / _Solutions.tex). Compiles with pdflatex.
% ======================================================================
\usepackage[utf8]{inputenc}
\usepackage[T1]{fontenc}
\usepackage[english]{babel}
\usepackage[a4paper,margin=2.4cm]{geometry}
\usepackage{amsmath,amssymb,amsthm,mathtools}
\usepackage{bm}
\usepackage{enumitem}
\usepackage{array}

\setlength{\parindent}{0pt}
\setlength{\parskip}{0.5em}
\emergencystretch=3em
\allowdisplaybreaks[1]

% ---------- math macros (same names as the rest of the project) ----------
\newcommand{\E}{\mathbb{E}}
\DeclareMathOperator{\Var}{Var}
\DeclareMathOperator{\Cov}{Cov}
\DeclareMathOperator{\Corr}{Corr}
\DeclareMathOperator*{\argmin}{arg\,min}
\DeclareMathOperator{\MSE}{MSE}
\DeclareMathOperator{\trace}{tr}
\newcommand{\Reals}{\mathbb{R}}
\newcommand{\Ints}{\mathbb{Z}}
\newcommand{\Comp}{\mathbb{C}}
\newcommand{\Nat}{\mathbb{N}}
\newcommand{\Prob}{\mathbb{P}}
\newcommand{\Tr}{^{\mathsf{T}}}
\newcommand{\conj}[1]{\overline{#1}}
\newcommand{\abs}[1]{\left\lvert #1 \right\rvert}
\newcommand{\norm}[1]{\left\lVert #1 \right\rVert}
\newcommand{\ind}{\mathbf{1}}
\newcommand{\eps}{\varepsilon}
\newcommand{\diff}{\nabla}
\newcommand{\acvs}{\gamma}
\newcommand{\acf}{\rho}
\newcommand{\pacf}{\alpha}
\newcommand{\sdf}{\mathcal{S}}
\newcommand{\Bsh}{B}
\newcommand{\iid}{\overset{\text{iid}}{\sim}}

\setlist[enumerate]{leftmargin=2em,topsep=2pt,itemsep=4pt}
\setlist[itemize]{leftmargin=1.6em,topsep=2pt,itemsep=2pt}

\newcounter{exo}

% ---------- exam cover header :  \examheader{exam title}{duration} ----------
\newcommand{\examheader}[2]{%
  \thispagestyle{empty}
  \begin{center}
    {\Large\bfseries Time Series}\\[3pt]
    Spring Semester 2026, EPFL\\
    \'Ecole Polytechnique F\'ed\'erale de Lausanne\\
    Prof.\ Dr.\ Sofia C.\ Olhede\\[7pt]
    \hrule height 0.8pt
    \vspace{9pt}
    {\large\bfseries Practice Exam --- #1}\\[4pt]
    {\normalsize Duration: #2 \quad$\cdot$\quad No calculator \quad$\cdot$\quad Answer on paper}\\
    \vspace{8pt}
    \hrule height 0.8pt
  \end{center}
  \vspace{14pt}
  \begin{tabular}{@{}p{8.2cm}@{}p{8cm}@{}}
    Name:\enspace\hrulefill & Surname:\enspace\hrulefill
  \end{tabular}\\[14pt]
  SCIPER (Student Card Number):\enspace\hrulefill
  \vspace{16pt}
}

% ---------- points table (fixed: 4 problems) ----------
\newcommand{\pointstable}{%
  \begin{center}
  \renewcommand{\arraystretch}{1.5}
  \begin{tabular}{|p{5cm}|p{4cm}|}
    \hline
    \textbf{Exercise} & \textbf{Points}\\\hline
    1 & \\\hline
    2 & \\\hline
    3 & \\\hline
    4 & \\\hline
    \textbf{Total} & \\\hline
  \end{tabular}
  \end{center}
}

% ---------- problem heading (exam: one problem per page) ----------
\newcommand{\problem}[1]{%
  \clearpage
  \stepcounter{exo}%
  \begin{center}\textbf{\large Problem \theexo\ (#1 points)}\end{center}
  \vspace{4pt}
}
% ---------- answer space: fill the statement page + one blank answer page ----------
% Put \answerpage at the END of each problem so every exercise gets its own page(s).
\newcommand{\answerpage}{\par\vfill\newpage\null\newpage}

% ---------- solutions document ----------
\newcommand{\solheader}[1]{%
  \begin{center}
    {\Large\bfseries Time Series --- Solutions}\\[3pt]
    {\large Practice Exam --- #1}\\[2pt]
    EPFL, MATH-342 \quad$\cdot$\quad Prof.\ S.\ C.\ Olhede\\[5pt]
    \hrule height 0.8pt
  \end{center}
  \vspace{10pt}
}
\newcommand{\solproblem}[1]{%
  \bigskip
  \stepcounter{exo}%
  {\large\bfseries Problem \theexo\ (#1 points)}\par\nobreak\vspace{2pt}
}
% Rappel de l'énoncé avant la solution (dans les corrigés)
\newcommand{\statementstart}{\par\vspace{1pt}{\bfseries\itshape Statement.}\par\nobreak\begingroup\small\itshape}
\newcommand{\statementend}{\par\endgroup\vspace{5pt}\hrule\vspace{6pt}{\bfseries Solution.}\par\nobreak\vspace{2pt}}

\begin{document}

\solheader{Week 6: Fourier Theory}

% ---------------------------------------------------------------
\solproblem{5}
\statementstart
Throughout this problem $g\in L^1$ has continuous-time Fourier transform
\[
G(f)=\int_{-\infty}^{\infty} g(t)\,e^{-2\pi i f t}\,dt ,
\]
and for a derived function $h$ we write $H$ for its transform. We use the four elementary Fourier-transform properties (conjugation, scaling, shift, linearity).

\textbf{(a)} Prove the \emph{conjugation} and \emph{time-shift} rules: if $h(t)=\conj{g(t)}$ then $H(f)=\conj{G(-f)}$, and if $h(t)=g(t+\tau)$ then $H(f)=G(f)\,e^{2\pi i f\tau}$. \hfill(2 pts)

\textbf{(b)} Combining scaling and shifting, show that the \emph{modulated} signal $h(t)=g(\alpha t)\,e^{2\pi i f_0 t}$ (with $\alpha\neq0$, $f_0\in\Reals$ fixed) has transform
\[
H(f)=\frac{1}{\abs\alpha}\,G\!\left(\frac{f-f_0}{\alpha}\right).
\]
State in one sentence what each of the two operations (scaling by $\alpha$, modulation by $e^{2\pi i f_0 t}$) does to the spectrum. \hfill(2 pts)

\textbf{(c)} Use part (a) to show that if $g$ is \emph{real and even} (so $g(t)=\conj{g(t)}=g(-t)$), then its transform $G$ is real and even. \hfill(1 pt)
\statementend

\textbf{(a)} \emph{Conjugation.} Pull the conjugate outside the integral; it flips the sign of the exponent:
\[
H(f)=\int_{-\infty}^{\infty}\conj{g(t)}\,e^{-2\pi i t f}\,dt
=\conj{\int_{-\infty}^{\infty} g(t)\,e^{2\pi i t f}\,dt}
=\conj{\int_{-\infty}^{\infty} g(t)\,e^{-2\pi i t(-f)}\,dt}
=\conj{G(-f)} .
\]
\emph{Time shift.} Insert $1=e^{-2\pi i f\tau}e^{2\pi i f\tau}$ to complete the exponent:
\[
H(f)=\int_{-\infty}^{\infty} g(t+\tau)\,e^{-2\pi i f t}\,dt
=\int_{-\infty}^{\infty} g(t+\tau)\,e^{-2\pi i f(t+\tau)}\,dt\;e^{2\pi i f\tau}
=G(f)\,e^{2\pi i f\tau},
\]
where in the last step we substituted $s=t+\tau$ ($ds=dt$, limits unchanged) so the integral is exactly $G(f)$.

\textbf{(b)} Build $h$ in two steps. First the \emph{scaling} rule: with $p(t)=g(\alpha t)$, substitute $s=\alpha t$ (so $dt=ds/\abs\alpha$, the $\abs\alpha$ arising from reversing the limits when $\alpha<0$),
\[
P(f)=\int_{-\infty}^{\infty} g(\alpha t)\,e^{-2\pi i f t}\,dt
=\frac{1}{\abs\alpha}\int_{-\infty}^{\infty} g(s)\,e^{-2\pi i (f/\alpha)s}\,ds
=\frac{1}{\abs\alpha}\,G\!\left(\frac{f}{\alpha}\right).
\]
Now $h(t)=p(t)\,e^{2\pi i f_0 t}$ is a pure \emph{modulation}; modulation in time is a frequency shift, since
\[
H(f)=\int_{-\infty}^{\infty} p(t)\,e^{2\pi i f_0 t}\,e^{-2\pi i f t}\,dt
=\int_{-\infty}^{\infty} p(t)\,e^{-2\pi i (f-f_0) t}\,dt
=P(f-f_0).
\]
Substituting $P$,
\[
\boxed{\,H(f)=P(f-f_0)=\frac{1}{\abs\alpha}\,G\!\left(\frac{f-f_0}{\alpha}\right).\,}
\]
In words: scaling time by $\alpha$ \emph{inversely} stretches the frequency axis (by $1/\alpha$) and rescales the height by $1/\abs\alpha$; modulating by $e^{2\pi i f_0 t}$ \emph{translates} the whole spectrum to be centred at $f_0$.

\textbf{(c)} Let $g$ be real and even. \emph{Real} ($\conj{g(t)}=g(t)$) means $g$ equals its own conjugate, so by the conjugation rule of (a),
\[
G(f)=\conj{G(-f)}\quad\Longleftrightarrow\quad G(-f)=\conj{G(f)}\qquad(\text{Hermitian symmetry}).
\]
\emph{Even} ($g(-t)=g(t)$) is the scaling rule of (b) with $\alpha=-1$ (so $\abs\alpha=1$), giving $G(-f)=G(f)$ ($G$ even). Combining the two,
\[
G(f)=G(-f)=\conj{G(f)},
\]
so $G(f)$ equals its own conjugate, i.e.\ $G$ is \textbf{real}; and $G(f)=G(-f)$ says $G$ is \textbf{even}. (This is the Fourier-theoretic reason a real, even signal always has a real, even transform.)

% ---------------------------------------------------------------
\solproblem{5}
\statementstart
Recall the continuous-time convolution of $g,h\in L^1$,
\[
(g*h)(t)=\int_{-\infty}^{\infty} g(t-y)\,h(y)\,dy ,
\]
and that upper-case letters denote Fourier transforms.

\textbf{(a)} Prove the \emph{convolution theorem}: if $u=g*h$ then $U=G\cdot H$. State explicitly the Fubini step and why the hypotheses $g,h\in L^1$ justify it. \hfill(3 pts)

\textbf{(b)} \emph{(Filtering.)} A signal $x\in L^1$ is passed through a linear filter with impulse response $a\in L^1$, producing $y=a*x$. Using part (a), give $Y$ in terms of $X$ and the transfer function $A$. If the filter is the running-average kernel $a(t)=\tfrac1{2T}\,\ind_{\{\abs t\le T\}}$, compute $A(f)$ in closed form and verify $A(0)=1$. \hfill(2 pts)
\statementend

\textbf{(a)} By definition of the transform and of convolution,
\[
\begin{aligned}
U(f)&=\int_{-\infty}^{\infty} u(x)\,e^{-2\pi i x f}\,dx
     =\int_{-\infty}^{\infty}\!\left(\int_{-\infty}^{\infty} g(x-y)\,h(y)\,dy\right)e^{-2\pi i x f}\,dx\\[2pt]
    &=\int_{-\infty}^{\infty}\!\left(\int_{-\infty}^{\infty} g(x-y)\,e^{-2\pi i x f}\,dx\right)h(y)\,dy
     \qquad\text{(Fubini: swap the order of integration)}\\[2pt]
    &=\int_{-\infty}^{\infty}\!\left(\int_{-\infty}^{\infty} g(s)\,e^{-2\pi i s f}\,ds\right)e^{-2\pi i y f}\,h(y)\,dy
     \qquad (s=x-y,\ dx=ds)\\[2pt]
    &=G(f)\int_{-\infty}^{\infty} h(y)\,e^{-2\pi i y f}\,dy=G(f)\,H(f).
\end{aligned}
\]
The Fubini interchange is legitimate because the double integral is absolutely convergent:
\[
\int_{-\infty}^{\infty}\!\!\int_{-\infty}^{\infty}\abs{g(x-y)\,h(y)}\,dx\,dy
=\int_{-\infty}^{\infty}\abs{g(s)}\,ds\int_{-\infty}^{\infty}\abs{h(y)}\,dy
=\norm g_1\,\norm h_1<\infty,
\]
using the change of variables $s=x-y$ and the hypotheses $g,h\in L^1$. Hence $U=G\cdot H$: the transform of a convolution is the product of the transforms.

\textbf{(b)} With $y=a*x$, part (a) gives immediately
\[
Y(f)=A(f)\,X(f),
\]
where $A(f)=\int_{-\infty}^{\infty} a(t)e^{-2\pi i f t}\,dt$ is the \emph{transfer function}: filtering multiplies the spectrum frequency-by-frequency by $A$. For the running average $a(t)=\tfrac1{2T}\ind_{\{\abs t\le T\}}$,
\[
A(f)=\frac{1}{2T}\int_{-T}^{T} e^{-2\pi i f t}\,dt
=\frac{1}{2T}\cdot\frac{e^{-2\pi i f t}}{-2\pi i f}\Bigg|_{t=-T}^{t=T}
=\frac{1}{2T}\cdot\frac{e^{2\pi i f T}-e^{-2\pi i f T}}{2\pi i f}
=\frac{\sin(2\pi f T)}{2\pi f T}.
\]
At $f=0$ the indeterminate form has limit $1$ (since $\sin x/x\to1$); equivalently $A(0)=\frac{1}{2T}\int_{-T}^{T}1\,dt=1$. So the filter preserves the zero-frequency (mean/DC) component and attenuates higher frequencies --- a low-pass filter, as a running average should be.

% ---------------------------------------------------------------
\solproblem{5}
\statementstart
This problem concerns the finite-data Fourier transform, the Dirichlet kernel, and aliasing. Take sampling interval $\Delta=1$ and observation times $t=0,1,\dots,n-1$.

\textbf{(a)} The transform of the finite rectangular window is the \textbf{Dirichlet kernel}
\[
D_n(f)=\sum_{t=0}^{n-1} e^{-2\pi i f t}.
\]
Sum this finite geometric series and show that
\[
D_n(f)=\frac{\sin(n\pi f)}{\sin(\pi f)}\;e^{-\pi i (n-1) f},
\qquad\text{so}\qquad
\abs{D_n(f)}=\left\lvert\frac{\sin(n\pi f)}{\sin(\pi f)}\right\rvert .
\]
\hfill(2 pts)

\textbf{(b)} Show that $\abs{D_n(f)}$ vanishes exactly at the nonzero \textbf{Fourier frequencies} $f_k=k/n$, $k=1,\dots,n-1$, while $D_n(0)=n$. (This is why the finite Fourier transform of a pure sinusoid sitting exactly on a Fourier frequency contributes nothing at the other Fourier frequencies.) \hfill(1.5 pts)

\textbf{(c)} \emph{(Aliasing / Nyquist.)} Let $g\in L^1$ have continuous-time Fourier transform $\mathcal G(f)$, and let $G$ denote the transform of the sampled sequence $\{g(t)\}_{t\in\Delta\Ints}$. Recall the aliasing identity
\[
G(f)=\sum_{k\in\Ints}\mathcal G\!\left(f+\frac{k}{\Delta}\right),
\qquad \abs f\le \frac{1}{2\Delta}.
\]
Suppose $\mathcal G$ is band-limited triangular,
$\mathcal G(f)=(1-\abs f)\,\ind_{\{\abs f\le 1\}}$.
For $\Delta=1$ (Nyquist frequency $\tfrac12$) compute the aliased transform $G(f)$ on $\abs f\le\tfrac12$, and state whether aliasing has occurred and why. \hfill(1.5 pts)
\statementend

\textbf{(a)} The sum is a finite geometric series with ratio $r=e^{-2\pi i f}$. For $f\notin\Ints$ (so $r\neq1$),
\[
D_n(f)=\sum_{t=0}^{n-1} r^{t}=\frac{1-r^{n}}{1-r}
=\frac{1-e^{-2\pi i n f}}{1-e^{-2\pi i f}} .
\]
Factor a half-angle phase out of numerator and denominator:
\[
\begin{aligned}
D_n(f)
&=\frac{e^{-\pi i n f}\bigl(e^{\pi i n f}-e^{-\pi i n f}\bigr)}
       {e^{-\pi i f}\bigl(e^{\pi i f}-e^{-\pi i f}\bigr)}
=\frac{e^{-\pi i n f}\,(2i)\sin(n\pi f)}{e^{-\pi i f}\,(2i)\sin(\pi f)}\\[2pt]
&=\frac{\sin(n\pi f)}{\sin(\pi f)}\;e^{-\pi i (n-1) f},
\end{aligned}
\]
using $e^{i\theta}-e^{-i\theta}=2i\sin\theta$. The phase factor has modulus $1$, so
\[
\abs{D_n(f)}=\left\lvert\frac{\sin(n\pi f)}{\sin(\pi f)}\right\rvert .
\]

\textbf{(b)} The numerator $\sin(n\pi f)$ vanishes exactly when $n f\in\Ints$, i.e.\ at $f=k/n$, $k\in\Ints$. The denominator $\sin(\pi f)$ vanishes only when $f\in\Ints$. At the Fourier frequencies $f_k=k/n$ for $k=1,\dots,n-1$ we have $f_k\notin\Ints$ (since $0<k/n<1$), so the denominator is nonzero while the numerator is zero:
\[
\abs{D_n(f_k)}=\frac{\abs{\sin(k\pi)}}{\abs{\sin(\pi k/n)}}=\frac{0}{\sin(\pi k/n)}=0,\qquad k=1,\dots,n-1.
\]
At $f=0$ both vanish; resolving the limit (or summing directly $\sum_{t=0}^{n-1}1$),
\[
D_n(0)=\lim_{f\to0}\frac{\sin(n\pi f)}{\sin(\pi f)}=\lim_{f\to0}\frac{n\pi f}{\pi f}=n .
\]
So $\abs{D_n}$ has a peak of height $n$ at the origin and is exactly zero at every \emph{other} Fourier frequency. Consequently the finite transform of a sinusoid sitting at one Fourier frequency contributes nothing at the remaining Fourier frequencies: there is no spectral leakage \emph{between} Fourier frequencies for an on-grid sinusoid.

\textbf{(c)} With $\Delta=1$ the Nyquist frequency is $\tfrac12$ and the aliasing identity reads $G(f)=\sum_{k\in\Ints}\mathcal G(f+k)$ on $\abs f\le\tfrac12$. Since $\mathcal G$ is supported on $[-1,1]$, for $\abs f\le\tfrac12$ only $k\in\{-1,0,1\}$ can contribute:
\[
G(f)=\mathcal G(f-1)+\mathcal G(f)+\mathcal G(f+1).
\]
Evaluate each term on $\abs f\le\tfrac12$:
\begin{itemize}
  \item $\mathcal G(f)=1-\abs f$ (always in support);
  \item $\mathcal G(f+1)=1-\abs{f+1}$ is nonzero only when $\abs{f+1}\le1$, i.e.\ $f\le0$, where it equals $1-(1+f)=-f$;
  \item $\mathcal G(f-1)=1-\abs{f-1}$ is nonzero only when $\abs{f-1}\le1$, i.e.\ $f\ge0$, where it equals $1-(1-f)=f$.
\end{itemize}
Hence
\[
G(f)=
\begin{cases}
(1-\abs f)+(-f)=(1+f)+(-f)=1, & -\tfrac12\le f\le 0,\\[2pt]
(1-\abs f)+f=(1-f)+f=1, & 0\le f\le \tfrac12,
\end{cases}
\]
that is, $\boxed{\,G(f)\equiv1\ \text{ on }\abs f\le\tfrac12.\,}$
\textbf{Aliasing has occurred.} The continuous transform $\mathcal G$ had support out to $\abs f=1$, i.e.\ \emph{beyond} the Nyquist frequency $\tfrac12$; sampling at $\Delta=1$ folds that excess high-frequency content back into the band $[-\tfrac12,\tfrac12]$. Here the folded triangular tails fill in exactly the missing area, turning a triangular continuous transform into a \emph{flat} sampled transform --- the high-frequency content has been irretrievably mixed into the low frequencies.

% ---------------------------------------------------------------
\solproblem{5}
\statementstart
\textbf{(Revision of Week 5: differencing \& SARIMA.)} Consider the signal-plus-noise model $X_t=\mu_t+Y_t$, where $\mu_t$ is a deterministic component and $\{Y_t\}$ is a zero-mean stationary process with autocovariance sequence $\acvs_\tau$. Recall $\diff=I-\Bsh$, $\diff_{(s)}=I-\Bsh^{s}$.

\textbf{(a)} Let $\mu_t=a+bt+s_t$ with a \emph{monthly} season $s_{t+12}=s_t$. Compute $\diff\,\diff_{(12)}X_t$ explicitly as a linear combination of the $Y$'s, verify its mean is $0$, and explain in one sentence why it is weakly stationary. \hfill(2 pts)

\textbf{(b)} Now let $\{X_t\}$ be the multiplicative seasonal process $\mathrm{SARMA}(0,1)\times(1,0)_{12}$,
\[
\bigl(1-\Phi\,\Bsh^{12}\bigr)X_t=\bigl(1-\theta\Bsh\bigr)\eps_t,
\qquad \abs\Phi<1 .
\]
By inverting the seasonal AR factor, obtain the $\mathrm{MA}(\infty)$ weights $\psi_j$ in $X_t=\sum_{j\ge0}\psi_j\eps_{t-j}$ (closed form for $\psi_{12l}$ and $\psi_{12l+1}$; all other $\psi_j=0$). \hfill(2 pts)

\textbf{(c)} Using $\acvs_\tau=\sigma^{2}\sum_{j\ge0}\psi_j\psi_{j+\tau}$, compute $\acvs_0$ and $\acvs_1$ in closed form, and evaluate $\acf_1=\acvs_1/\acvs_0$ for $\theta=\tfrac12$. \hfill(1 pt)
\statementend

\textbf{(a)} Apply the seasonal difference first. Using $s_t=s_{t-12}$,
\[
\diff_{(12)}X_t=X_t-X_{t-12}
=\bigl(a+bt+s_t+Y_t\bigr)-\bigl(a+b(t-12)+s_{t-12}+Y_{t-12}\bigr)
=12b+\bigl(Y_t-Y_{t-12}\bigr).
\]
The period-$12$ season is annihilated and the linear trend collapses to the constant $12b$. Now apply $\diff=I-\Bsh$, which kills the constant:
\[
\diff\,\diff_{(12)}X_t
=\bigl(12b+Y_t-Y_{t-12}\bigr)-\bigl(12b+Y_{t-1}-Y_{t-13}\bigr)
=\boxed{\,Y_t-Y_{t-1}-Y_{t-12}+Y_{t-13}\,}.
\]
This is a \emph{fixed} finite linear combination of $\{Y_t\}$, so its mean is
$\E[Y_t-Y_{t-1}-Y_{t-12}+Y_{t-13}]=0$, and being a fixed linear filter of a stationary process its autocovariance depends on the lag $\tau$ only (no $t$); hence it is weakly stationary. All $t$-dependence (trend and season) has been removed.

\textbf{(b)} Solve the model for $X_t$: $X_t=\Phi X_{t-12}+\eps_t-\theta\eps_{t-1}$, a seasonal $\mathrm{AR}(1)$ at lag $12$ driven by an $\mathrm{MA}(1)$ input. Invert the seasonal factor as a geometric series in $\Bsh^{12}$ (valid since $\abs\Phi<1$),
\[
\frac{1}{1-\Phi\Bsh^{12}}=\sum_{l\ge0}\Phi^{l}\Bsh^{12l},
\]
so that
\[
X_t=\sum_{l\ge0}\Phi^{l}\Bsh^{12l}\bigl(\eps_t-\theta\eps_{t-1}\bigr)
=\sum_{l\ge0}\Phi^{l}\bigl(\eps_{t-12l}-\theta\,\eps_{t-12l-1}\bigr).
\]
Reading off the coefficient of $\eps_{t-j}$,
\[
\boxed{\;\psi_{12l}=\Phi^{l},\qquad \psi_{12l+1}=-\theta\,\Phi^{l}\quad(l\ge0),\qquad \psi_j=0\ \text{otherwise.}\;}
\]
(Check: $\psi_0=1,\ \psi_1=-\theta,\ \psi_{12}=\Phi,\ \psi_{13}=-\theta\Phi,\dots$)

\textbf{(c)} Use $\acvs_\tau=\sigma^{2}\sum_{j\ge0}\psi_j\psi_{j+\tau}$. The only nonzero taps are at $j\equiv0$ and $j\equiv1\pmod{12}$.

\emph{Lag $0$:}
\[
\acvs_0=\sigma^{2}\sum_{l\ge0}\bigl(\psi_{12l}^{2}+\psi_{12l+1}^{2}\bigr)
=\sigma^{2}\sum_{l\ge0}\bigl(\Phi^{2l}+\theta^{2}\Phi^{2l}\bigr)
=\frac{(1+\theta^{2})\,\sigma^{2}}{1-\Phi^{2}} .
\]
\emph{Lag $1$:} the only surviving pairs are $\psi_{12l}\,\psi_{12l+1}=\Phi^{l}(-\theta\Phi^{l})=-\theta\Phi^{2l}$, so
\[
\acvs_1=\sigma^{2}\sum_{l\ge0}\bigl(-\theta\Phi^{2l}\bigr)=\frac{-\theta\,\sigma^{2}}{1-\Phi^{2}} .
\]
Therefore
\[
\acf_1=\frac{\acvs_1}{\acvs_0}=\frac{-\theta}{1+\theta^{2}} .
\]
For $\theta=\tfrac12$: $\acf_1=\dfrac{-1/2}{1+1/4}=\dfrac{-1/2}{5/4}=-\dfrac{2}{5}=-0.4.$ (The autocorrelation is governed entirely by the $\mathrm{MA}(1)$ side-factor near the origin; the seasonal $\mathrm{AR}(1)$ produces analogous spikes at the seasonal lags $12,24,\dots$ of relative size $\Phi^{k}$.)

\end{document}
